%%%%%
%
%% Dans ce fichier, les �ventuels accents sont cod�s
%% en ISO-8859-15 (aka latin9)
%% La ligne suivante donne trois exemples 
%% e accent aigu : �, e accent grave : �, e-dans-l'o : �
\section{Module d'un vecteur en 2 dimensions}

\itemstyle{Op�rateur~:} \operfmt{Module\_2D\_flt}
\index{Module\_2D\_flt}

%%%%%OBJET%%%%
\itemstyle{Objet~:}
Calcule le module (norme euclidienne) d'une 
quantit� vectorielle � deux dimensions (typiquement une
vitesse, un vent horizontal) dont les composantes
sont contenues dans deux vecteurs (au sens AIDA)
distincts.

%%%%%ARGUMENTS%%%%
\bigskip
%% Dans ce fichier, les �ventuels accents sont cod�s
%% en ISO-8859-15 (aka latin9)
%% La ligne suivante donne trois exemples 
%% e accent aigu : �, e accent grave : �, e-dans-l'o : �
%%%%%ENTREES%%%%
\begin{tabarg}
\multicolumn{3}{l}{\titretab {Entr\'ees}}\\
\hline
C$_1$ & Vecteur, \texttt{float} & Premi�re composante du vecteur (unit� quelconque)\\
C$_2$ & Vecteur, \texttt{float} & Deuxi�me composante du vecteur (m�me unit� que C$_1$)\\
Name & Coeff., \texttt{string} & Description courte du r�sultat \\
Eng\_name & Coeff., \texttt{string} & Description courte du r�sultat en langue anglaise \\
\end{tabarg}

\medskip
%%%%%SORTIES%%%%
\begin{tabarg}
\multicolumn{3}{l}{\titretab {Sorties}}\\
\hline
M & Vecteur, \texttt{float} & Module du vecteur (m�me unit� que les 2 composantes)\\
\end{tabarg}


%%%%%FORMULE%%%%
\itemstyle{Formule}

%% Formule de l'operateur : Module_2D_flt
%% Dans ce fichier, les �ventuels accents sont cod�s
%% en ISO-8859-15 (aka latin9)
%% La ligne suivante donne trois exemples 
%% e accent aigu : �, e accent grave : �, e-dans-l'o : �
\[
M = \sqrt{ C_1^2 + C_2^2 }
\]



%%%%%COMMENTAIRE%%%%
\itemstyle{Commentaire~:}
Cet op�rateur prend en compte les vecteurs $C_{1}$ et $C_{2}$ 1D et 2D, tant que $C_{1}$ et $C_{2}$ sont de la m�me dimension.

%%%%%AUTEUR%%%%
\itemstyle{Auteur~:}
CNRM/GMEI/TRAMM

%%%%%REFERENCE%%%%
\itemstyle{R�f�rences~:}
Sans objet
